\section{Discussion}
\label{sec:slam-discussion}


\subsection{Baseline GLIM Performance Analysis}

The evaluation of the baseline GLIM algorithm revealed two critical failure modes that fundamentally compromised its localisation reliability: systemic vertical drift and a catastrophic ``point of no return'' in planar estimation.

\subsubsection{Vertical Drift and Altitude Recovery}
The baseline algorithm exhibited a pronounced systemic drift in the positive altitude ($z$-axis) direction. This phenomenon is clearly illustrated in Dataset 4 (Figure~3.5g), where the vehicle traverses a linear trajectory down a vineyard row. Although the actual environment features a modest elevation gain of approximately 5~m, the baseline GLIM algorithm accumulated a significant vertical error, culminating in nearly 30~m of drift by the end of the row.

Interestingly, this drift demonstrated a self-correcting property during the return traversal. As the vehicle moved in the opposite direction, the vertical error decreased at a rate symmetrical to its initial accumulation, eventually reconverging with the ground truth data. This behavior suggests that the initial drift was effectively encoded into the generated LiDAR map. During the return pass, the scan-matching process was constrained by this previously mapped geometry. Consequently, the odometry "undid" the accumulated drift by aligning the current scans with the (erroneous) elevation of the existing map. 

Notably, this vertical distortion of the map did not appear to hinder subsequent scan matching or loop closure. The algorithm remained capable of registering scans against a map with significant Z-axis distortion. This resilience may be a byproduct of the specific experimental geometry—characterized by repetitive A-B linear paths with minimal lateral movement—which allows the scan matcher to maintain a high correlation in the X-Y plane despite vertical inconsistencies.

\subsubsection{Planar Failure and the 'Point of No Return'}
In contrast to the recoverable nature of vertical drift, errors in the X-Y plane frequently led to catastrophic and irrecoverable failure. This "point of no return" is most evident in Figure~3.4a, where the baseline algorithm maintains high fidelity for the first half of the traversal before experiencing a sudden, exponential increase in error. 

At this critical threshold, the accumulated planar drift distorts the LiDAR map to an extent that subsequent scan matching becomes fundamentally unreliable. Once in this failure state, the trajectory becomes increasingly jagged and stochastic, likely caused by false-positive scan matches that introduce large, discontinuous jumps in the pose estimation. 

This failure mode is particularly prevalent in the structured and repetitive environments of vineyard rows. Small angular or lateral drifts cause parallel rows in the LiDAR map to diverge or overlap incorrectly. Once this spatial distortion exceeds the scan matcher's convergence radius, the algorithm may incorrectly register the vehicle's current position to an adjacent row. This ``row-skipping'' effect creates a feedback loop of map corruption, where each incorrect match further distorts the global map, rendering recovery unlikely for the remainder of the trajectory.

\subsection{Refined GLIM Performance}

The transition from the baseline GLIM algorithm to the refined (parameterised) version demonstrates a significant enhancement in localisation stability and drift suppression. While the baseline algorithm frequently succumbed to the failure modes described in previous sections, the refined version exhibits a more controlled error growth profile.

\subsubsection{Magnitude of Error Reduction}
The most prominent improvement is observed in the \textit{error percent per metre} metric. In the baseline configuration, Dataset 2 exhibited a peak error of 7.3610\%, whereas the refined version reduced this to 3.6150\%, a reduction of approximately 50\%. In more favourable conditions, such as Dataset 4, the improvement was even more dramatic: the error dropped from 4.1910\% in the baseline to a significantly lower 0.6680\% in the refined version. This indicates that parameter optimisation effectively mitigates the ``point of no return'' by maintaining better geometric alignment over longer durations.

\subsubsection{Relative Pose Stability}
The \textit{RPE 1 m translation standard deviation} provides insight into the ``smoothness'' or ``jitter'' of the trajectory. While the baseline GLIM suffered from jagged trajectories around points of failure -- peaking at a standard deviation of 0.9384 in Dataset 3 -- the refined GLIM consistently lowered this jitter across all datasets. For instance, in Dataset 3, the refined version achieved an RPE standard deviation of 0.8239. Although this remains the highest in the refined category, the overall reduction suggests that the optimised parameters facilitate more consistent scan-matching, reducing the frequency of the ``jumps'' caused by false-positive matches in repetitive environments.

\subsubsection{Drift Suppression per Meter}
The \textit{error per metre (m)} metric further reinforces these findings. The refined GLIM consistently achieved sub-centimetre drift per metre in datasets where the baseline was losing several centimetres (e.g., Dataset 0 dropping from 0.0340~m to 0.0087~m). This suggests that the refinement process successfully tuned the fusion weights and scan-matching thresholds to better suit the specific LiDAR characteristics and environment structure, delaying the onset of the catastrophic failures observed in the baseline.

\begin{table}[H]
\centering
\caption{Comparative Performance: Baseline vs. Refined GLIM}
\label{tab:comparison_summary}
\begin{tabularx}{\textwidth}{lXXXXX}
\toprule
\textbf{Metric} & \textbf{Dataset 0} & \textbf{Dataset 1} & \textbf{Dataset 2} & \textbf{Dataset 3} & \textbf{Dataset 4} \\ \midrule
\textbf{Error \% Reduction} & 73.9\% & 51.4\% & 50.8\% & 84.6\% & 84.0\% \\
\textbf{RPE Std Improvement} & 12.2\% & 5.8\% & 28.8\% & 12.2\% & 19.9\% \\ \bottomrule
\end{tabularx}
\end{table}

\subsection{Proposed Method: GNSS-Augmented Localisation}

The integration of GNSS measurements into the GLIM factor graph constitutes a fundamental architectural change: by introducing a source of absolute positional observability absent from both the baseline and refined GLIM configurations, it eliminates the unbounded drift inherent in scan-matching-only odometry. The resulting improvement of approximately two orders of magnitude in localisation accuracy relative to the no-GNSS baseline is consistent with this design intent, validating the fusion architecture rather than representing an unexpected outcome. The proposed method's error percentage per metre ranges from $0.0153\%$ (Dataset 4) to $0.0501\%$ (Dataset 2), compared to the refined GLIM's range of $0.6081\%$ to $3.6150\%$. This reduction is consistent across all datasets, with no observable catastrophic failure modes analogous to those seen in earlier iterations.

\subsubsection{GNSS Integration and Global Constraint}

The performance improvement stems fundamentally from the incorporation of GNSS-derived positional constraints within the factor graph optimisation. Whereas GLIM operates exclusively on scan-matching odometry, the proposed method augments this local relative information with periodic absolute position measurements from the GNSS receiver. This integration serves two critical functions: (1) it provides global observability that prevents unbounded drift accumulation, and (2) it allows the factor graph optimiser to ``correct'' local odometric errors retroactively across the entire trajectory window.

The effectiveness of this approach is evident in comparing error accumulation profiles. In the refined GLIM, even with parameter optimisation, the error percent per metre remains linear or superlinear across the traversal. By contrast, the proposed method maintains sub-centimetre per-metre errors throughout, suggesting that the global constraints are sufficiently frequent and accurate to prevent drift from establishing itself in the first place.

This result is further reinforced by the GNSS factor utilisation statistics (Table~3.7 in the results section): across all datasets, between 0.65--0.78\% of all available GNSS datapoints were injected into the factor graph following outlier rejection. This sparsity is itself a key finding of this work. A primary motivation for the proposed architecture was to reduce operational dependence on RTK-GNSS, which is susceptible to canopy occlusion, multipath interference, weather conditions, and infrastructure cost in agricultural deployments. The results demonstrate that the vast majority of localisation accuracy is recovered through the LiDAR-inertial odometry backbone, with GNSS contributing only infrequent global anchors to prevent long-term drift accumulation. Achieving two orders of magnitude improvement over the no-GNSS baseline at sub-percent GNSS utilisation confirms that dense or continuous RTK coverage is not a prerequisite for precise trajectory estimation in this architecture.

\subsubsection{Dataset-Specific Performance Variability}

Despite the strong performance of the proposed method across all datasets, performance does vary across datasets in a pattern that reveals environmental factors influencing localisation quality. Dataset 2 exhibits the highest error metrics (0.0501\% error per metre, RMSE $y$ of 0.2003~m, ATE of 0.1741~m), whereas Datasets 1 and 3 achieve the lowest errors (0.0154\% error per metre). This variability is instructive rather than concerning, as the absolute error magnitudes remain within acceptable bounds for agricultural manipulation.

The degradation in Dataset 2 is likely attributable to environmental conditions rather than algorithmic limitations. Visual inspection of the traversal pattern and GNSS utilisation (0.65\% of available GNSS datapoints injected into the factor graph, the lowest across all datasets) suggests that portions of the trajectory experienced reduced GNSS signal quality or multipath interference. The geometric consistency-based outlier rejection, whilst effective at filtering genuinely erroneous points, cannot recover lost signal availability. However, even under these challenging conditions, the localisation error remained controlled at $0.0501\%$, indicating robustness to transient GNSS degradation.

\subsubsection{Relative Pose Error and Short-Horizon Consistency}

The Relative Pose Error (RPE) metrics, which assess the consistency of local pose increments, provide insight into frame-to-frame localisation smoothness. The proposed method achieves RPE 1~m translation standard deviations ranging from $0.0205$ to $0.0479$~m across all datasets, compared to the refined GLIM's $0.2449$ to $0.8239$~m. This two-to-tenfold reduction in local pose jitter indicates that the global constraints are not merely correcting large-scale drift, but are improving the consistency of short-horizon relative measurements.

This improvement likely arises from the factor graph's ability to reason about odometric uncertainty jointly with global constraints. When the optimiser encounters a region where local odometry is inconsistent with global GNSS observations, it propagates the correction smoothly across the pose graph rather than applying isolated jump corrections. This smoothness is essential for downstream applications such as robotic arm control, where abrupt discontinuities in perceived position could cause trajectory tracking errors.

\subsubsection{GASP Metric Behaviour and Optimisation Dynamics}

A notable feature of the results is a slight upward trend in the GASP (Global Alignment and Smoothness Profile) metric for the single best-performing run across the four parameter optimisation stages. Unlike the baseline GLIM error percentage, which exhibits clear convergence and monotonic improvement, the proposed method's best-run GASP metric increases marginally from stage one to stage four. This trend is attributable to the stochastic nature of the Latin Hypercube parameter search rather than to any systematic degradation: the differences between near-optimal configurations are negligible, and stage one already identifies a configuration that fully exploits the GNSS integration architecture. Subsequent stages sample nearby regions of the parameter space and, by chance, return configurations that score marginally worse on GASP, without producing any material change in trajectory accuracy. The magnitude of the variation is small enough that it falls within the expected noise of the pseudo-random search process.

Despite this marginal decline in peak performance, the overall distribution of GASP metric values improves significantly across the four stages. As visible in \cref{fig:proposed_error_dist}, the error distributions shift progressively leftward over successive stages, indicating an improvement in typical performance. Concurrently, the distributions become taller and narrower — that is, the standard deviation decreases — reflecting a marked increase in the reliability and consistency of results. In contrast, the stage zero distribution is comparatively flat and wide, indicating that results at that stage were highly variable and unpredictable. This evolution demonstrates that the optimisation process is not merely identifying a single lucky configuration, but is systematically converging the parameter space towards a regime in which good performance is reproducible. The trade-off between peak and median performance is therefore a natural consequence of this convergence: as the search focuses on consistently performant regions, extreme but unreliable high-performing outliers are replaced by a tighter cluster of reliably good solutions.

This observation underscores an important principle: no single metric fully captures the multidimensional nature of SLAM performance. GASP is most informative when comparing qualitatively different configurations (e.g., baseline vs. GNSS-augmented); within a near-optimal neighbourhood, trajectory accuracy metrics such as ATE and RPE are more reliable discriminators.

\subsection{Comparative Analysis and Mechanisms of Improvement}

Examining the three approaches in concert reveals a clear hierarchy of performance and elucidates the mechanisms driving improvements at each stage.

\subsubsection{Parametrisation vs. Sensor Fusion}

The transition from baseline to refined GLIM demonstrates that parameter tuning can reduce error magnitude significantly (50--85\% reduction in error percentage). However, parameter optimisation alone does not eliminate the fundamental instability of scan-matching-only odometry in repetitive vineyard environments. The refined GLIM still exhibits occasional trajectory inconsistencies and elevated RPE metrics relative to the proposed method.

By contrast, the transition from refined GLIM to the GNSS-augmented method represents a qualitative shift in localisation architecture. Rather than optimising the tuning parameters of a single sensor modality, the proposed method introduces a fundamentally independent source of global position information. This architectural change achieves improvements that parametrisation alone cannot reach. The error reduction from refined GLIM to the proposed method ($\sim$50--98\%) is larger and more consistent than the improvement from baseline to refined ($\sim$50--85\%), and the proposed method exhibits no dataset in which performance degrades significantly.

\subsubsection{Role of Outlier Rejection}

The geometric consistency-based outlier rejection procedure is a critical enabling factor in this result. The procedure admits only those GNSS measurements that satisfy pairwise distance-preservation constraints across the candidate window, rejecting the large majority of available raw measurements before factor graph integration. This aggressive filtering is precisely what makes the sparse utilisation viable: by ensuring that every incorporated constraint is geometrically consistent, the procedure prevents a single erroneous global measurement from corrupting the factor graph. The architecture therefore achieves the desired property of GNSS-frugal localisation — using the global sensor as a rare but reliable anchor rather than a continuous dependency. This deterministic filtering approach, based on rigid-body distance preservation, is well-suited to the repetitive vineyard environment, where probabilistic alternatives such as RANSAC would require outlier distribution assumptions that may not hold in structured agricultural settings.

The robustness of this design is confirmed by Dataset 2, where GNSS factor utilisation falls to its lowest observed value of 0.65\%, yet localisation error remains controlled at $0.0501\%$ per metre. This indicates that each geometrically verified constraint carries sufficient positional information to anchor the trajectory even when GNSS availability is degraded.

\subsubsection{Geometric Quality Assessment and Kabsch Stability}

The planar geometry quality assessment phase, which enforces minimum requirements on translational baseline, spread ratio, rotational excitation, and condition number, ensures that the Kabsch algorithm operates in a numerically stable regime. This preprocessing step prevents ill-conditioned covariance matrices that could yield biased or unstable alignment solutions.

The quintet of measurement from Datasets 0--4 all satisfy the geometry quality thresholds, which explains the consistency of global alignment across the evaluation. In datasets with weakly excited geometry (e.g., highly linear traversals with minimal rotation), this step would conservatively defer global alignment until sufficient geometric excitation is accumulated, preventing premature or inaccurate initialisation.

\subsection{Implications for Vineyard Robotics}

The proposed method's consistent sub-centimetre accuracy has direct implications for robotic pruning feasibility. Pruning tasks require localisation accuracy sufficient to position the robotic arm within reach of target vine structures. With typical vine stems ranging from 5--15~mm in diameter and clusters requiring workspace accessibility, the proposed method's maximum absolute trajectory error of approximately 0.8--1.0~m (observed across all datasets) represents an acceptable error bound. This error magnitude can be accommodated through local perception-based arm guidance or online trajectory adjustment based on close-range stereo or depth sensing.

Furthermore, the low local pose jitter (RPE 1~m standard deviation of $0.0205$--$0.0479$~m) is important for smooth arm trajectory execution. Abrupt jumps in perceived position, typical of the baseline and refined GLIM approaches, would cause tracking servo errors and potentially jerky or unsafe arm motion. The proposed method's smoothness mitigates this risk.

However, a critical limitation remains unaddressed by the current evaluation: all experiments were conducted at slow traversal speeds (approximately 0.5~m/s) with ample GNSS signal availability. Real vineyard deployments will encounter variable canopy occlusion, faster navigation speeds during inter-row transit, and periods of GNSS dropout beneath dense foliage. The robustness of the method under these challenging real-world conditions has not yet been empirically demonstrated. Chapter 4 will address this limitation by evaluating precise alignment techniques that can operate with minimal or absent GNSS input, using instead local dense perception and pre-existing 3D models to constrain pose estimation.